\section{6. Discussion}
\subsection{6.1 Why stability is the correct narrative frame}

Three a-priori plausible framings for the contributions of this paper exist:


\begin{enumerate}
\item \textbf{Accuracy lift on single-tool selection} (Subtask1). Observed $\Delta \le +6$pp under strict label-logprob scorers; $+11$pp under legacy substring (scorer-dependent).
\item \textbf{Accuracy lift on multi-tool selection} (Subtask4 F-simultaneous). Originally predicted $+5$–$+15$pp; empirically FALSIFIED at full 497 ($-4.6$pp).
\item \textbf{Direction-specificity of the ontology subspace} (null-control gaps). Observed $+16$ to $+49$pp on Subtask1 (scorer-invariant); $+68.5$pp on Subtask4 (full 497, verified at the largest scale in the paper).
\end{enumerate}

Only framing (3) is robustly supported at the magnitude predicted by theory (Cor 6.9.6). The $\pm 5$pp headlines of framings (1) and (2) are scorer-dependent and task-specific; the $+30$–$+68$pp gaps of (3) are scorer-invariant, task-invariant, and cross-model (§5.4 Table). We therefore lead with the stability claim. Accuracy lifts, when they occur (Qwen sum +0.10 / mean +5.03, Llama-Base sum +6.33 / mean +2.61, Mistral-Base sum +3.12, MMLU flat $\alpha=0.2$ +1.4, contrastive Subtask4 smoke +5.8), are supporting evidence that the direction is \emph{downstream-usable}, not the main contribution.


This framing also restructures the paper's falsifiability. Under the accuracy-lift narrative, the paper has a single failure point (Subtask4 $-4.6$pp already observed). Under the stability narrative, the main claim is already verified at full scale; accuracy-lift extensions (§5.5.2 contrastive, §5.10.1 LoRA) are independent follow-ups whose individual success or failure leaves the main contribution intact.


\subsection{6.1.1 Why the cross-model positive is supporting (not leading) evidence}

Qwen + Llama-Base + Mistral-Base sum-positive triad is \textbf{generalization evidence for the stability claim}: the ontology direction is uniquely privileged in three independent transformer families, not a Qwen-specific artifact. The Mistral-Instruct-v0.3 negative is not a counterexample to stability — its no\_steer itself is 7.84pp below Mistral-Base (61.51\% vs 69.35\%), and the further $-2.92$pp shift under $\alpha=0.3$ K-bias is consistent with chat-template hedging rather than a failure of the ontology direction to be privileged. A null-control comparison on Mistral-Instruct (random/featshuffle at $\alpha=0.3$) is queued and predicted to show the same $+60$+pp direction-specificity gap as Qwen, confirming stability universality with Instruct-family hedging as a separate scope limit.


\subsection{6.2 (R) as a design constraint, not a technicality}

Section 3.2's hard-gate MMLU degradation is not a bug — it is the direct empirical signal that regularity matters. Design the gate to be Lipschitz; do not select a hard threshold for nominal interpretability.


\subsection{6.3 Why K-side, not Q-side}

[AdaSEKA differentiation per \texttt{adaseka\_vs\_ours\_differentiation\_2026\_04\_10.md}]. Q-side 1-of-M routing is structurally capped at rank \texttt{r} (Cor 6.9). K-side F-simultaneous attains rank \texttt{R}. For multi-facet intents the gap becomes operationally relevant on compositional benchmarks (Sec 5.11).


\subsection{6.4 Limitations}

\begin{enumerate}
\item Qwen + Llama only. Mistral requires base-weakness mitigation (Instruct variant, Sec 5.3 H2).
\item Generation scorer + label\textbackslash\{\}\_logprob disagree at N=20 smoke; full 995 pending.
\item Compositional benchmark is the highest leverage axis; BFCL-v3 integration deferred to Wave 4.
\end{enumerate}
